\documentclass[tikz,border=6pt]{standalone}
\usepackage{ctex}
\usepackage{amsmath}
\usetikzlibrary{calc, arrows.meta, decorations.markings}

\begin{document}
\begin{tikzpicture}[scale=1.0, thick,
  every node/.style={font=\small}]

%% -----------------------------------------------
%% 公理 1：两点确定一条直线（左侧子图）
%% -----------------------------------------------
\begin{scope}[xshift=0cm]
  % 直线
  \draw[black, thick] (-1.2, 0) -- (2.8, 0);
  % 两点 A、B
  \coordinate (A1) at (0, 0);
  \coordinate (B1) at (2, 0);
  \fill[black] (A1) circle [radius=2pt];
  \fill[black] (B1) circle [radius=2pt];
  \node[below] at (A1) {$A$};
  \node[below] at (B1) {$B$};
  % 直线标签
  \node[above right] at (2.4, 0) {$l$};
  % 标题
  \node[align=center, font=\footnotesize] at (0.8, -1.0)
    {公理 1：过两点有且只有\\一条直线 $AB \subset l$};
\end{scope}

%% -----------------------------------------------
%% 公理 2：三点确定一个平面（中间子图）
%% -----------------------------------------------
\begin{scope}[xshift=4.5cm]
  % 平面（平行四边形表示）
  \coordinate (PL) at (-1.2, -0.2);
  \coordinate (PR) at (2.3, -0.2);
  \coordinate (PRt) at (2.8, 0.8);
  \coordinate (PLt) at (-0.7, 0.8);
  \fill[gray!20] (PL) -- (PR) -- (PRt) -- (PLt) -- cycle;
  \draw[gray] (PL) -- (PR) -- (PRt) -- (PLt) -- cycle;
  % 三点（不共线）
  \coordinate (A2) at (0.0, 0.1);
  \coordinate (B2) at (1.5, 0.05);
  \coordinate (C2) at (0.7, 0.65);
  \fill[black] (A2) circle [radius=2pt];
  \fill[black] (B2) circle [radius=2pt];
  \fill[black] (C2) circle [radius=2pt];
  \node[below left] at (A2) {$A$};
  \node[below right] at (B2) {$B$};
  \node[above] at (C2) {$C$};
  % 辅助线（三点构成三角形）
  \draw[red, dashed] (A2) -- (B2) -- (C2) -- cycle;
  % 标题
  \node[align=center, font=\footnotesize] at (0.8, -1.0)
    {公理 2：不共线三点\\确定唯一平面 $\alpha$};
\end{scope}

%% -----------------------------------------------
%% 公理 3：两平面公共点构成直线（右侧子图）
%% -----------------------------------------------
\begin{scope}[xshift=9.5cm]
  % 平面 alpha（稍左偏）
  \coordinate (aLL) at (-1.0, -0.5);
  \coordinate (aLR) at (0.9, -0.5);
  \coordinate (aUR) at (1.3, 0.6);
  \coordinate (aUL) at (-0.6, 0.6);
  \fill[blue!10] (aLL) -- (aLR) -- (aUR) -- (aUL) -- cycle;
  \draw[blue!50] (aLL) -- (aLR) -- (aUR) -- (aUL) -- cycle;
  \node[blue!70, font=\footnotesize] at (-0.8, 0.5) {$\alpha$};

  % 平面 beta（稍右偏，穿插 alpha）
  \coordinate (bLL) at (-0.2, -0.6);
  \coordinate (bLR) at (2.0, -0.3);
  \coordinate (bUR) at (1.8, 0.7);
  \coordinate (bUL) at (-0.4, 0.4);
  \fill[red!10] (bLL) -- (bLR) -- (bUR) -- (bUL) -- cycle;
  \draw[red!50] (bLL) -- (bLR) -- (bUR) -- (bUL) -- cycle;
  \node[red!70, font=\footnotesize] at (1.7, 0.6) {$\beta$};

  % 交线 l（两平面的公共直线）
  \draw[black, very thick] (0.0, -0.45) -- (1.0, 0.55);
  \node[right, font=\footnotesize] at (1.0, 0.55) {$l$};
  \fill[black] (0.0, -0.45) circle [radius=2pt];
  \fill[black] (1.0, 0.55) circle [radius=2pt];

  % 标题
  \node[align=center, font=\footnotesize] at (0.8, -1.0)
    {公理 3：两平面有公共点\\则交于一条直线 $\alpha \cap \beta = l$};
\end{scope}

\end{tikzpicture}
\end{document}
